\section{Related Work}
\label{sec:related}

Algorithmic collusion has been a long-standing concern in oligopoly economics and is now a first-order regulatory question.\footnote{Concurrent benchmark work \citep{Tailor_2025_whisper} introduces a multi-domain collusion-detection benchmark spanning several economic and adversarial settings; our scope is the basin-of-attraction characterization at fixed Bertrand-pricing setting and is complementary in scope.} \citet{ezrachi2016virtual} introduce a four-category taxonomy --- Messenger, Hub-and-Spoke, Predictable Agent, Digital Eye --- and argue that the Predictable Agent category, in which self-learning algorithms converge independently to supra-competitive outcomes without explicit coordination, is the case existing antitrust frameworks are least equipped to detect: U.S.\ and EU competition law typically requires evidence of explicit agreement, which self-learning agents do not produce. \citet{asker2021ai_pricing} reframe the question productively: ``does AI collude'' is the wrong frame, because collusion is a property of the learning rule and game structure, not of AI per se. The right question is which algorithmic design choices induce collusion, and whether learning-rule classes generalize. Our paper sits inside this frame: we extend the empirical case beyond the $Q$-learner family to LLM agents, with the basin-of-attraction structure as the unit of comparison.

\paragraph{$Q$-learner basins as the empirical baseline.}
\citet{calvano_etal_2020_ai_algorithmic_pricing_collusion} establish the canonical Bertrand-pricing $Q$-learner result: two independent $Q$-learning agents converge to supra-competitive prices in repeated Bertrand duopoly, with terminal price levels distributed across multiple distinct attractors as a function of the initial $Q$-matrix and exploration parameters. The Calvano supplementary appendix histograms over 1000+ runs are the closest prior analog of our cross-rep terminal-state distribution. \citet{waltman_kaymak_2008_qlearning_cournot} report a parallel result in repeated Cournot oligopoly --- $Q$-learning agents converge to terminal quantities ranging from Cournot-Nash to collusive monopoly across independent replications, with the distribution of terminal quantities multimodal --- demonstrating that equilibrium selection under reinforcement learning is the rule, not the exception, in algorithmic oligopoly games. \citet{wunder_littman_babes_2010_qlearning_basin} formalize the basin-of-attraction lens for $Q$-learning: terminal-state multimodality is the empirical signature of basin structure, and the basin a trajectory lands in depends on initial conditions, exploration schedule, and the topology of the value-function landscape. These three works establish multimodal terminal distributions as the recognized signature of basin structure under reinforcement-learning-based pricing; our LLM result is the natural cross-family question.

\paragraph{LLM agents as a third learning family.}
LLMs are a qualitatively distinct class of agents from $Q$-learners: they carry no shared reward signal, no explicit value function, no online weight updates, no exploration parameter. Their ``learning'' within a pricing game is in-context, mediated through textual reasoning and prompt history, not by gradient updates to a value table. \citet{akata2023playing} report early evidence that LLMs play repeated games in non-trivial ways and exhibit cooperative or strategic behaviors not predicted by their training objective alone. \citet{wei2022cot} establish chain-of-thought reasoning as an emergent capability with measurable behavioral effects on multi-step tasks. \citet{brown2024monkeys} document broad emergent agency in tool-using LLMs. More directly relevant: \citet{Cao_2026_llm_collusion} report bistable pricing-game outcomes for LLM agents in a Bertrand-style setting, anticipating the bimodal structure we recover at GATE-5; and \citet{Hammond_2025_multiagent_risks,Riedl_2025_emergent} survey emergent multi-agent LLM behaviors and the systemic risks they create at scale. Our paper supplies the focused, pre-registered, multi-condition basin-of-attraction characterization that this literature has not yet produced for the Bertrand setting.

\paragraph{Empirical antitrust evidence.}
The phenomenon is field-relevant, not just simulation-relevant. \citet{assad_etal_2024_algorithmic_pricing_competition_empirical} document supra-competitive algorithmic pricing in retail markets with empirical-IO methods; \citet{fish2024algorithmicCollusion} extend the case with field experiments demonstrating collusive convergence in production deployments. \citet{BrownMackay_2024_nber_pricing_cadence} show that pricing-cadence asymmetries between competing algorithms generate the same collusive outcomes Calvano predicts in simulation, providing a cousin signal that asymmetric dynamic structure --- not just symmetric multimodality --- drives collusion in algorithmic pricing. Our convergence-rate asymmetry between $n = 2$ and $n = 5$ is consistent with the Brown-Mackay finding that asymmetric dynamics shape collusive outcomes; we extend it cross-condition rather than cross-firm.

\paragraph{Equilibrium-selection and basin theory.}
The theoretical apparatus for characterizing basins predates the algorithmic-collusion literature. \citet{harsanyi_selten_1988_general_theory_equilibrium_selection} introduce risk-dominance and payoff-dominance as criteria for selecting among multiple Nash equilibria. \citet{young_1993_evolution_of_conventions} formalizes equilibrium selection under stochastic learning dynamics; \citet{fudenberg_levine_1998_learning_in_games} survey learning-in-games convergence results across reinforcement, fictitious play, and best-response classes. \citet{benaim_hofbauer_sorin_2005_stochastic_approximation_differential_inclusions} provide the stochastic-approximation foundation for analyzing differential inclusions that govern basin structure under noisy learning. \citet{sandholm_2010_population_games} extend the framework to population dynamics. \citet{askenazi2024bounds} derive Folk-Theorem-style bounds on the set of supra-competitive prices achievable by $Q$-learning and gradient-learning agents in repeated pricing games, and show that the basin of attraction for collusive equilibria depends on initialization and exploration rate --- linking algorithmic collusion to the equilibrium-selection literature in evolutionary game theory. The basin-width-by-agent-count hypothesis we develop in \S\ref{sec:discussion} extends this framework empirically to the LLM pricing setting; we treat agent count as a structural game property that modulates effective in-context exploration, with the formal $Q$-learner bounds of \citeauthor{askenazi2024bounds} as a theoretical reference rather than a direct prediction for LLM agents.

\paragraph{Gap and positioning.}
Two specific gaps the prior literature does not close. First, whether the LLM family of agents reproduces the basin-of-attraction signature observed in $Q$-learners has not been characterized at multiple agent counts under a pre-registered, multi-rep design. The closest precedents are \citet{Cao_2026_llm_collusion} (LLM bistability without the cross-condition basin-width comparison) and \citet{Tailor_2025_whisper} (multi-domain benchmark without the per-condition basin characterization). Second, whether basin width depends on agent count, and whether smaller $n$ yields wider basins, has been theoretically predicted \citep{askenazi2024bounds,wunder_littman_babes_2010_qlearning_basin} but not directly tested in any LLM pricing study. Unlike \citet{calvano_etal_2020_ai_algorithmic_pricing_collusion} (Q-learning, two agents only), \citet{Cao_2026_llm_collusion} (LLM bistability without cross-condition contrast), and \citet{Tailor_2025_whisper} (multi-domain benchmark without per-condition basin characterization), our approach is a pre-registered focused study of basin structure under one fixed Bertrand setting at two distinct agent counts, with hierarchical integrity scaffolding (cryptographic pre-registration $+$ cross-toolchain verification $+$ adversarial review) and a Stage 4 cross-model replication path explicitly committed.
